\section{極形式と複素平面}\label{節：極形式と複素平面}
\subsection{複素数を点として表す}

$i^2=-1$を満たす虚数単位$i$を用い,\,複素数を$z=x+iy$と書く.
$x,y$は実数である.
$x$を実部,$y$を虚部といい,\,平面上の点$(x,y)$と対応させる.
横軸を実軸,\,縦軸を虚軸と呼ぶこの平面が\textbf{複素平面}である.

複素数の和は成分ごとの和だから,\,平面の矢印の和と一致する.
積は$i^2=-1$を使って計算する.
例えば
\[
 i(x+iy)=-y+ix
\]
であり,\,点$(x,y)$は$(-y,x)$へ移る.
これは原点のまわりの反時計回りの直角回転である.

$z=x+iy$に対して$\overline z=x-iy$を\textbf{共役複素数}という.
実軸に関する鏡映に当たり,
\[
 z\overline z=x^2+y^2=|z|^2
\]
となる. $|z|=\sqrt{x^2+y^2}$を絶対値という.
$z\ne0$なら
\[
 \frac1z=\frac{\overline z}{|z|^2}
\]
だから,\,複素数でも零以外で割ることができる.

\subsection{極形式と積}

零でない$z$を,\,原点からの距離$\rho>0$と角度$\theta$で表す.
\[
 z=\rho(\cos\theta+i\sin\theta).
\]
これを\textbf{極形式}という.
角度$\theta$は偏角と呼ばれ,$2\pi$の整数倍を加えても同じ点を表す.
零には偏角を定めない.

加法公式
\[
 \begin{aligned}
 \cos(\theta+\phi)&=\cos\theta\cos\phi-\sin\theta\sin\phi,\\
 \sin(\theta+\phi)&=\sin\theta\cos\phi+\cos\theta\sin\phi
 \end{aligned}
\]
を使うと,
\[
 \begin{aligned}
 &(\cos\theta+i\sin\theta)(\cos\phi+i\sin\phi)\\
 &\qquad=\cos(\theta+\phi)+i\sin(\theta+\phi).
 \end{aligned}
\]
従って積では距離を掛け,\,角度を足す.
商では距離を割り,\,角度を引く.

$\lambda=\rho(\cos\theta+i\sin\theta)$を固定して掛けると,\,原点からの距離は$\rho$倍,\,角度は$\theta$だけ進む.
$\rho=1$なら一定角の回転,$\rho<1$なら縮小を伴う回転となる.
\NewFigRotation

積の公式を繰り返せば,\,非負整数$m$について
\[
 (\cos\theta+i\sin\theta)^m
 =\cos(m\theta)+i\sin(m\theta).
\]
これが\textbf{ド・モアブルの定理}である.
本書では$\cos\theta+i\sin\theta$を$e^{i\theta}$とも略記する.
ここでは回転を表す記号として使い,\,複素指数関数の一般論は必要としない.

\subsection{商は長さの比と角度の差}

二点$\alpha,\beta$を基準にして,\,点$z$の位置を
\[
 w=\frac{z-\alpha}{z-\beta}\qquad(z\ne\beta)
\]
で記録してみよう.
絶対値は
\[
 |w|=\frac{|z-\alpha|}{|z-\beta|}
\]
だから,\,二つの基準点からの距離の比を表す.
偏角は,\,ベクトル$z-\beta$から$z-\alpha$への有向角である.
二つの矢印を同じ始点に平行移動して角度を比較してもよい.
\NewFigRatio

この比が一定倍率で変わるなら,\,距離の比と角度の差の変化を同時に追える.
次節で,1次分数型に現れる置き換えがこの形になる.

\begin{bookpoint}
複素数の積は\textbf{倍率と回転角}をまとめて記録する.
複素数の商は\textbf{長さの比と角度の差}をまとめて記録する.
\end{bookpoint}
